% ==================================================================
%  PostProcessing.tex  --  Section 3.3 "Client-side Post-processing"
%  Loaded from main.tex.
%  Implementation (Client_main.py, assign_clusters_by_gap): noise band
%  0.5 label unit; gap threshold = one label unit, estimated as
%  max(value) / (largest labelled slot + 1).
% ==================================================================

\subsection{Client-side Post-processing}
\label{sec:post}

The client decrypts the ciphertext it receives and recovers one value per
slot. These values are real numbers rather than exact multiples of the
label unit $1/n$. Each maximum is approximate
(Appendix~\ref{app:error}), and every CKKS operation adds a small noise,
so the labels drift slightly over the rounds.

The client therefore separates the clusters by the gaps between the label
values rather than by rounding them. The points of one cluster carry the
same label up to these small errors, so their values lie close together.
The labels of two clusters started at different multiples of $1/n$, so
their values lie about one label unit or more apart. Points that received
no label from a core point keep a value near $0$ and are reported as
noise.
